% Загальний аналіз опрацьованої літератури до теми
% «Моделювання потоку плазми в контурі токамаку»
% Компіляція: xelatex (двічі) або latexmk -xelatex
\documentclass[11pt,a4paper]{article}

\usepackage{fontspec}
\usepackage{polyglossia}
\setdefaultlanguage{ukrainian}
\setotherlanguage{english}
\setmainfont{cmunrm}[Extension=.otf,BoldFont=cmunbx,ItalicFont=cmunti,BoldItalicFont=cmunbi]
\setsansfont{cmunss}[Extension=.otf,BoldFont=cmunsx,ItalicFont=cmunsi,BoldItalicFont=cmunso]
\setmonofont{cmuntt}[Extension=.otf,BoldFont=cmuntb,ItalicFont=cmunit,BoldItalicFont=cmuntx]

\usepackage{amsmath,amssymb}
\usepackage[margin=2.2cm]{geometry}
\usepackage{booktabs,longtable,array,tabularx}
\usepackage{enumitem}
\usepackage[table]{xcolor}
\usepackage[colorlinks=true,linkcolor=blue!50!black,citecolor=green!40!black,urlcolor=blue!60!black]{hyperref}
\usepackage{xurl}

\setlist{itemsep=2pt,topsep=4pt}
\newcommand{\rel}[1]{\textbf{#1}}
\newcommand{\file}[1]{\texttt{\small #1}}
\renewcommand{\arraystretch}{1.15}

\title{Моделювання потоку плазми в контурі токамаку:\\
       загальний аналіз опрацьованої літератури}
\author{Огляд корпусу джерел із теки \file{library}}
\date{28 вересня 2026 р.}

\begin{document}
\maketitle

\begin{abstract}
Проаналізовано 32 повнотекстові документи (1139 сторінок), завантажені з відкритих
джерел за списком літератури до теми «Моделювання потоку плазми в контурі токамаку»
(ключові слова: JET, динаміка плазми, токамак). Документи класифіковано за вісьмома
тематичними кластерами й оцінено за релевантністю темі за шкалою 0–5. Ядро корпусу
утворюють роботи з неокласичної теорії та експериментальних даних про тороїдальні й
полоїдальні потоки плазми і радіальне електричне поле $E_r$
\cite{Singh2004,Sigmar1987,Seto2014,Ida2001}. На їх основі запропоновано каркас
математичної моделі для проєкту. Головна прогалина корпусу --- відсутність робіт,
що безпосередньо моделюють потоки на JET: JET фігурує лише як верифікаційний розряд
або довідкові параметри \cite{Hoppe2022,Tyshchenko2021,vanVugt2019}.
\end{abstract}

\tableofcontents

%======================================================================
\section{Мета, обсяг і методика аналізу}

\paragraph{Мета.} Визначити, які з опрацьованих джерел можуть лягти в основу
математичної моделі потоку плазми в токамаці, які --- дають контекст або дані для
валідації, а які є непрофільними; виявити прогалини в корпусі.

\paragraph{Обсяг корпусу.} Зі списку літератури (\file{«Моделювання потоку плазми в
контурі токамаку».md}, $\approx 80$ позицій) у відкритому доступі вдалося отримати
32 документи:
\begin{itemize}
  \item CORE (агрегатор відкритого доступу) --- 20 документів;
  \item Google Scholar (дисертації ХФТІ, ХНУ, УФЖ) --- 5;
  \item Наукова періодика України (НБУВ) --- 3;
  \item Репозитарій НТУ «ХПІ» --- 4.
\end{itemize}
Не отримано: ELAKPI (13 позицій; репозитарій доступний лише з мережі КПІ), друковані
видання фонду Бібліотеки КПІ (6), статті SpringerLink (12; з них 4 у відкритому
доступі, але сайт блокує автоматичне завантаження), три класичні статті
\emph{Nuclear Fusion} (Арцимович 1972, Furth 1975, Муховатов--Шафранов 1971), книга
R.~B.~White, а також матеріали з файлообмінника twirpx (свідомо не завантажувались).

\paragraph{Методика.}
\begin{enumerate}
  \item Витяг текстового шару (\texttt{pdftotext}); для сканів 1987--1992 рр. і PDF
        зі зламаним кодуванням шрифтів зміст і формули звірено з зображеннями сторінок.
  \item Транскрипція всіх PDF у Markdown із формулами LaTeX інструментом
        \textbf{Marker} (datalab, v2.0) --- результати в \file{library/transcripts/}.
  \item Для кожного документа складено структуроване резюме: мета, модель і методи,
        ключові рівняння, результати, релевантність (0--5), обмеження.
  \item Зведення: класифікація, синтез моделі, аналіз прогалин.
\end{enumerate}

\paragraph{Шкала релевантності.} 5 --- безпосередня основа моделі потоку;
4 --- пряма фізика потоків або дані для валідації; 3 --- суттєвий методичний
внесок (коди, транспорт, JET-дані); 2 --- корисний контекст; 1 --- периферійно;
0 --- непридатне/ненадійне.

%======================================================================
\section{Загальна характеристика корпусу}

\begin{table}[h]
\centering
\caption{Розподіл документів за тематичними кластерами}
\label{tab:clusters}
\begin{tabularx}{\textwidth}{@{}lXc@{}}
\toprule
№ & Кластер & К-сть \\
\midrule
1 & Рівновага, МГД-стійкість, керування положенням і формою & 5 \\
2 & Транспорт, обертання і потоки плазми (неокласика, турбулентність, домішки) & 9 \\
3 & Крайова плазма, дивертор, взаємодія плазма--стінка, пил & 5 \\
4 & Кінетика: електрони-утікачі, альфвенові хвилі, частинки в полі & 4 \\
5 & Інтегроване моделювання і коди & 3 \\
6 & Інженерія та матеріали установки & 2 \\
7 & Загальні / оглядові & 2 \\
8 & Непрофільні & 2 \\
\midrule
 & \textbf{Разом} & \textbf{32} \\
\bottomrule
\end{tabularx}
\end{table}

\paragraph{Хронологія.} 1980-ті --- 3; 1990-ті --- 4; 2000-ні --- 8;
2010-ті --- 11; 2020-ті --- 6. Класичні неокласичні роботи (MIT, PPPL, 1987--1992)
співіснують із сучасними інтегрованими кодами (Fenix 2022, STREAM 2022).

\paragraph{Мови.} Англійська --- 20, українська --- 6, російська --- 6.
Усі англомовні документи --- з CORE; вітчизняні джерела здебільшого
стосуються суміжних тем (механіка котушок, пилова плазма, розрядна плазма).

\paragraph{Релевантність.} 5 балів --- 2 документи; 4 --- 2; 3 --- 7; 2 --- 7;
1 --- 12; 0 --- 2. Отже, лише $\approx 1/3$ корпусу (11 документів із балом $\ge 3$)
безпосередньо працює на тему; решта --- контекст.

\begin{longtable}{@{}p{0.9cm}>{\raggedright\arraybackslash}p{4.3cm}p{0.6cm}p{0.5cm}>{\raggedright\arraybackslash}p{8.6cm}@{}}
\caption{Зведена таблиця джерел (Кл.\ --- кластер, Р --- релевантність 0--5)}
\label{tab:all}\\
\toprule
Код & Автори, рік & Кл. & Р & Суть \\
\midrule
\endfirsthead
\toprule
Код & Автори, рік & Кл. & Р & Суть \\
\midrule
\endhead
\bottomrule
\endfoot
C01 & Singh, Rogister, Kaw, 2004 \cite{Singh2004} & 2 & \rel{5} &
  Неокласичне тороїдальне обертання на краю при полоїдально асиметричному напуску
  газу; 1D-рівняння кутового моменту; зв'язок з $E_r$ і L--H переходом. \\
C02 & Zou, Li, Lv, 2013 \cite{Zou2013} & 1 & 1 &
  Теореми існування для задачі Ґреда--Шафранова з вільною межею та нелокальним тиском. \\
C03 & Sigmar, Zanino, Hsu, 1987 \cite{Sigmar1987} & 2 & \rel{4} &
  Рідинна модель паралельних/полоїдальних/тороїдальних потоків, $E_r$ і полоїдальної
  асиметрії домішок в обертовій плазмі; рівняння для $\omega\propto E_r$. \\
C04 & Wong, Cheng, 1987 \cite{Wong1987} & 2 & 3 &
  Обертання підсилює неокласичну дифузію важких домішок на 1--2 порядки (TFTR). \\
C05 & Ward, Jardin, 1992 \cite{Ward1992} & 1 & 1 &
  Деформованість плазми і стабілізація вертикальної моди (NOVA-W). \\
C06 & Hoppe, Ekmark, Berger, Fülöp, 2022 \cite{Hoppe2022} & 4 & 1 &
  0D-код STREAM: утікачі й burn-through на старті; верифікація на JET \#77210. \\
C07 & Mazzucato, 2000 \cite{Mazzucato2000} & 7 & 1 &
  Скейлінгова оцінка ролі $B_T$ для експерименту з палаючою плазмою. \\
C08 & Fable et al., 2022 \cite{Fable2022} & 5 & 3 &
  «Тренажер польоту» Fenix (ASTRA + SPIDER): повний розряд ASDEX Upgrade, EU-DEMO. \\
C09 & Park, Goldston, Taylor, 1989 \cite{Park1989} & 2 & 2 &
  Інверсія інтерферометрії TFTR; $\tau_E$ корелює з піковістю густини. \\
C10 & Seto, Fukuyama, 2014 \cite{Seto2014} & 2 & \rel{5} &
  Єдина 2D-багаторідинна транспортна модель для ядра і SOL з неокласичним
  замиканням; основа коду TASK/T2. \\
C11 & van Vugt, Huijsmans et al., 2019 \cite{vanVugt2019} & 2 & 3 &
  JOREK + трасування W у повній орбіті: $E\times B$-перенос під час ELM (ASDEX Upgrade). \\
C12 & Ida et al., 2001 \cite{Ida2001} & 2 & \rel{4} &
  Експеримент: структура $E_r$, спонтанні тороїдальні/полоїдальні потоки
  (JFT-2M, CHS, LHD); недифузійний потік імпульсу $\propto\nabla T_i$. \\
C13 & Hanada et al., 1990 \cite{Hanada1990} & 1 & 1 &
  Придушення пилкоподібних коливань локалізованим ЕЦН (WT-3, PEST). \\
C14 & Aho-Mantila, 2011 \cite{AhoMantila2011} & 3 & 3 &
  SOLPS/ERO/ASCOT проти експериментів ASDEX Upgrade: дрейфи визначальні, потоки в
  SOL занижено у 2--3 рази. \\
C15 & Pigarov et al., 2006 \cite{Pigarov2006} & 3 & 1 &
  Перенос пилу (DUSTT/UEDGE), DIII-D, NSTX. \\
C16 & Merlo, 2011 \cite{Merlo2011} & 1 & 2 &
  Вільнограничний розв'язувач Ґреда--Шафранова, реконструкція межі, LQG-керування
  формою (RFX-mod). \\
C17 & Hulse, 1992 \cite{Hulse1992} & 2 & 2 &
  1D-модель зарядових станів домішок, дифузія + пінч; чутливість до атомних даних. \\
C18 & Jardin et al., 2000 \cite{Jardin2000} & 5 & 2 &
  Бенчмарк коду TSC на NSTX; балонна межа $\beta\approx12\,\%$. \\
C19 & Fukuyama, Kasai, Furutani, 1992 \cite{Fukuyama1992} & 5 & 2 &
  1.5D-код TASK/TR для палаючої плазми класу ITER. \\
C20 & Wang, Lin, \dots, Lee et al., 2006 \cite{Wang2006} & 2 & 3 &
  Гірокінетичний код GTC у загальній геометрії: ITG-турбулентність і зональні потоки. \\
G1 & Марценюк, 2025 \cite{Martsenyuk2025} & 3 & 1 &
  Тліючі та комбіновані розряди для кондиціонування стінок (Ураган-2М); 0D-модель. \\
G2 & Марчук, 2025 \cite{Marchuk2025} & 4 & 2 &
  0D-модель утікачів при зриві (T-10, EAST), НГР-хвилі, ермітова інтерполяція
  швидкісних коефіцієнтів. \\
G3 & Ромащенко, 2021 \cite{Romashchenko2021} & 3 & 1 &
  Заряджання та нагрів макрочастинок у пучково-плазмових системах (автореферат). \\
G4 & Тищенко, 2021 \cite{Tyshchenko2021} & 4 & 3 &
  Альфвенові хвилі, острови, каналювання енергії $\alpha$-частинок; 1D-баланс
  $T_e,T_i$ для розрядів JET DTE1. \\
G5 & Толочкевич та ін., 2015 \cite{Tolochkevych2015} & 8 & 1 &
  PIC-моделювання кільватерних полів згустків у незамагніченій плазмі. \\
Н1 & Аксенов, 2015 \cite{Aksenov2015} & 8 & 0 &
  Нестандартна теорія «самозбудження» поля в токамаці; науково ненадійна. \\
Н2 & Дуба, 2018 \cite{Duba2018} & 1 & 2 &
  Simulink-модель керування положенням шнура (3 входи, 3 виходи). \\
Н3 & Москвитина, 2010 \cite{Moskvitina2010} & 4 & 3 &
  Повноорбітне трасування $\alpha$-частинки в ITER з гофруванням поля; Рунге--Кутта 5-го порядку. \\
Х1 & Зайцев, Асаёнок, 2003 \cite{Zaitsev2003} & 6 & 1 &
  Термонапруження гвинтової обмотки торсатрона (МСЕ). \\
Х2 & Зайцев, Шульженко, 2001 \cite{Zaitsev2001} & 6 & 1 &
  Керування напруженим станом котушки тороїдального поля. \\
Х3 & Зимников, 2016 \cite{Zimnikov2016} & 7 & 0 &
  Студентські тези про термоядерний синтез (популярний рівень). \\
Х4 & Беляева та ін., 2016 \cite{Belyaeva2016} & 3 & 1 &
  Рекристалізація W-дзеркал ITER під нейтронами і розпиленням. \\
\end{longtable}

%======================================================================
\section{Тематичний аналіз}

\subsection{Потоки, обертання плазми та радіальне електричне поле (ядро теми)}

Це центральний кластер для проєкту. Чотири документи
\cite{Singh2004,Sigmar1987,Seto2014,Ida2001} разом дають повний ланцюжок
«теорія $\to$ рівняння $\to$ експеримент».

\paragraph{Радіальний баланс сил.} Спільним елементом усіх робіт кластера є
зв'язок радіального електричного поля з потоками і градієнтом тиску іонів:
\begin{equation}\label{eq:Er}
  E_r=\frac{1}{Z_i e n_i}\frac{\partial p_i}{\partial r}-v_{\theta i}B_\varphi+v_{\varphi i}B_\theta .
\end{equation}
Singh та ін.\ \cite{Singh2004} використовують (\ref{eq:Er}) для пояснення впливу
напуску газу на поріг L--H переходу; Ida \cite{Ida2001} наголошує, що в геліотронах
$E_r$ визначається переважно полоїдальним потоком ($E_r\approx v_\theta B_\varphi$);
Seto і Fukuyama \cite{Seto2014} отримують (\ref{eq:Er}) як рівняння нульового порядку
за ларморівським радіусом у довільних координатах магнітних поверхонь.

\paragraph{Структура потоку на магнітній поверхні.} Sigmar, Zanino, Hsu
\cite{Sigmar1987} розкладають швидкість на паралельну компоненту і тверде обертання
поверхні з частотою $\omega(\psi)\propto E_r$:
\begin{equation}\label{eq:V0}
  \mathbf V^{(0)}=\Big(V_\parallel+\omega\frac{I}{B}\Big)\mathbf b-\omega R^2\nabla\varphi,
  \qquad \omega=c\,\frac{\partial\Phi}{\partial\psi},\quad I=RB_\varphi .
\end{equation}
Усереднене по поверхні рівняння для $\omega$ виявляється квадратичним,
\begin{equation}\label{eq:omega}
  \frac{\partial\omega}{\partial t}+C_1\omega^2+C_2\omega+C_3=0,
\end{equation}
де $C_1$ --- інерція, $C_2$ --- гальмування і паралельна в'язкість, $C_3$ ---
джерела імпульсу (NBI). Два корені означають можливу біфуркацію стаціонарного
обертання --- важливий якісний результат для чисельного моделювання.

\paragraph{Рівняння тороїдального імпульсу.} Singh та ін.\ \cite{Singh2004}
отримують 1D-рівняння дифузійного типу для густини кутового моменту з
перезарядним гальмуванням, турбулентним/неокласичним радіальним потоком імпульсу
і неокласичним джерелом, що залежить від полоїдальної асиметрії нейтралів
$N_n=N_{0n}(1+\Delta_1\cos\chi)$:
\begin{equation}\label{eq:mom}
  \frac{\partial}{\partial t}\big(m_iN_iU_{\varphi}\big)+\nu_{cx}m_iN_iU_{\varphi}
  +\frac{1}{r}\frac{\partial}{\partial r}\Big[r\,m_iN_i\big(\langle\tilde u_E\tilde u_{\varphi}\rangle+U_rU_{\varphi}\big)\Big]
  =S_{\rm neo}\big(\partial_rU_\parallel,\partial_r\ln T_i,\Delta_1\big).
\end{equation}
Оцінка $U_\varphi\sim0.2\,c_i\sim10$~км/с без асиметрії та до $\sim30$~км/с при
напуску з внутрішнього боку є, за зауваженням нашого аналізу, радше якісною: оцінений
у тій самій роботі $\Delta_1\sim10^{-3}$ значно менший за значення, використане для
ілюстрації.

\paragraph{Експериментальні орієнтири.} Ida \cite{Ida2001} (JFT-2M, CHS, LHD)
показує, що потік тороїдального імпульсу містить недифузійний доданок:
\begin{equation}\label{eq:Gamma}
  \Gamma_\varphi=-\chi_\varphi\frac{\partial(m_in_iv_\varphi)}{\partial r}+\alpha\,\nabla T_i,
  \qquad \alpha/\chi_\varphi\sim0.1\text{--}0.3,
\end{equation}
причому напрям спонтанного (intrinsic) обертання в токамаках і геліотронах
протилежний і визначається симетрією поля. Числові дані ($E_r$ до $+15$~кВ/м,
полоїдальні потоки 10--20~км/с, тороїдальні до $5\times10^4$~м/с) придатні як
порядки величин для валідації.

\paragraph{Двовимірне узагальнення.} Seto і Fukuyama \cite{Seto2014} формулюють
замкнену систему (8 рівнянь переносу на сорт + 5 електромагнітних) у координатах
магнітних поверхонь $(\rho,\chi,\zeta)$, коректну і всередині, і поза сепаратрисою.
Паралельна в'язкість зшиває форму Хіршмана--Сігмара (ядро) з Брагінським (SOL),
а $E_\rho$ визначається із закону Гаусса, а не з квазінейтральності. Після
усереднення модель відтворює неокласичне полоїдальне обертання. Робота суто
формальна (без чисельних результатів), але найкраще підходить як теоретична основа.

\paragraph{Вплив обертання на транспорт.} Wong і Cheng \cite{Wong1987} показують,
що при $m_I\omega^2R^2\gg T_I$ неокласична дифузія важких домішок зростає на
1--2 порядки (у режимі Пфірша--Шлютера $D_{\rm neo}$ залежить лише від $\omega^2$ і
$\omega^4$). van Vugt та ін.\ \cite{vanVugt2019} і Wang та ін.\ \cite{Wang2006}
демонструють домінантну роль $E\times B$-потоку --- відповідно під час ELM
(10--20\,\% екранованого W проникає всередину) та в регуляції ITG-турбулентності
зональними потоками.

\paragraph{Висновок до кластера.} Усі теоретичні роботи обмежені режимом
Пфірша--Шлютера (висока зіткненнєвість), великим аспектним відношенням і коловим
перерізом. Для гарячої плазми JET (бананний режим у ядрі) потрібні відповідні
неокласичні коефіцієнти, яких корпус не містить.

\subsection{Транспорт і інтегроване моделювання розряду}

Роботи \cite{Fukuyama1992,Fable2022,Jardin2000,Hulse1992,Park1989,Tyshchenko2021}
задають стандартну 1.5D-архітектуру «транспорт + рівновага + кола котушок».

Базова система в стандартній формі (TASK/TR \cite{Fukuyama1992}; аналогічні рівняння в ASTRA
\cite{Fable2022}, TSC \cite{Jardin2000}, у моделі Тищенка \cite{Tyshchenko2021}):
\begin{align}
  \frac{\partial n_s}{\partial t}&=-\frac1r\frac{\partial}{\partial r}(r\Gamma_s)+S_s,
  &\Gamma_s&=v_sn_s-D_s\frac{\partial n_s}{\partial r},\label{eq:n}\\
  \frac{\partial}{\partial t}\Big(\tfrac32n_sT_s\Big)&=-\frac1r\frac{\partial}{\partial r}(rQ_s)+P_s,
  &Q_s&=-n_s\chi_s\frac{\partial T_s}{\partial r}+\tfrac52T_s\Gamma_s,\label{eq:T}\\
  \frac{\partial B_\theta}{\partial t}&=\frac{\partial E_z}{\partial r},
  &E_z&=\eta\Big(\frac{1}{\mu_0r}\frac{\partial(rB_\theta)}{\partial r}-J_{\rm NB}-J_{\rm BS}\Big).\label{eq:B}
\end{align}
Замикання $\chi_s$ варіюється від моделей дрейфово-хвильової турбулентності
з довжиною змішування \cite{Fukuyama1992} та гіро-Бомівської основи з неокласикою
\cite{Fable2022} до сталих значень, узятих із TRANSP для JET
($\chi_i=5000$~см$^2$/с, $\chi_e=2700$~см$^2$/с) \cite{Tyshchenko2021}.

\paragraph{Цінність для проєкту.}
\begin{itemize}
  \item Тищенко \cite{Tyshchenko2021} --- єдина робота корпусу з власними
        розрахунками для JET: 1D-баланс $T_e,T_i$ для розрядів DTE1 (\#41069,
        \#42856, \#42847) з NBI, $\alpha$-нагрівом і каналюванням енергії хвилями.
        Результат: $T_{i0}$ зростає з 13.35 (NBI) до 15.71 (NBI+$\alpha$) і 16.15~кеВ
        (каналювання, $\nu=0.3$). Модель проста й відтворювана --- хороший перший тест.
  \item Fable та ін.\ \cite{Fable2022} --- приклад валідації часових трас повного
        розряду проти експерименту; обертання і перенос імпульсу, однак, не моделюються.
  \item Hulse \cite{Hulse1992} --- неявна схема для $Z+1$ зв'язаних рівнянь
        зарядових станів; наголошує, що похибка атомних даних удвічі еквівалентна
        похибці $D$ удвічі.
\end{itemize}
Жоден документ кластера не містить рівняння тороїдального імпульсу в
1.5D-системі --- його доведеться додати на основі (\ref{eq:mom})--(\ref{eq:Gamma}).

\subsection{Рівновага, МГД-стійкість і керування положенням}

Рівняння Ґреда--Шафранова є спільною основою \cite{Merlo2011,Zou2013,Jardin2000,Fable2022}:
\begin{equation}\label{eq:GS}
  \Delta^*\psi\equiv R\frac{\partial}{\partial R}\Big(\frac1R\frac{\partial\psi}{\partial R}\Big)
  +\frac{\partial^2\psi}{\partial Z^2}=-\mu_0R^2\frac{dp}{d\psi}-F\frac{dF}{d\psi}.
\end{equation}
Merlo \cite{Merlo2011} дає практичний розв'язувач із вільною межею (Maxfea),
параметризацію $J_\varphi(\bar\psi)$ і методи реконструкції межі (згадано XLOC, що
використовується на JET), а також лінеаризовану модель плазма--провідники
$L^*\dot{\mathbf I}+R\mathbf I=\mathbf V$ з LQG-регулятором. Ward і Jardin
\cite{Ward1992} та Дуба \cite{Duba2018} стосуються вертикальної стабілізації;
модель Дуби (кожен канал --- $1/(s^2+k)$) надто спрощена і не обґрунтована.
Zou та ін.\ \cite{Zou2013} --- чиста математика (існування розв'язків).
Hanada та ін.\ \cite{Hanada1990} --- придушення пилкоподібних коливань.

Для проєкту кластер дає \emph{геометрію}: сітку магнітних поверхонь, на якій
розв'язуються рівняння потоку, та зв'язок транспорт~$\leftrightarrow$~рівновага
(схема FCT у \cite{Seto2014}).

\subsection{Крайова плазма, дивертор, стінка}

Aho-Mantila \cite{AhoMantila2011} --- найцінніше джерело кластера: валідація
SOLPS5.0/ERO/ASCOT на ASDEX Upgrade. Ключові висновки для моделі потоку:
без дрейфів $E\times B$ і діамагнітного розв'язки SOL незадовільні; паралельні потоки
в головному SOL модель занижує у 2--3 рази (пряме поле). Граничні умови
на мішені:
\begin{equation}\label{eq:Bohm}
  M=\frac{u_\parallel}{c_s}\ge1,\qquad c_s=\sqrt{\frac{ZT_e+\gamma T_i}{m_i}},\qquad
  \frac{eV_{sh}}{T_e}\approx0.5\ln\Big[2\pi\frac{m_e}{m_i}\Big(1+\frac{T_i}{T_e}\Big)\Big].
\end{equation}
Pigarov та ін.\ \cite{Pigarov2006}, Ромащенко \cite{Romashchenko2021},
Марценюк \cite{Martsenyuk2025}, Беляева та ін.\ \cite{Belyaeva2016} стосуються пилу,
шарів біля стінки, кондиціонування та ерозії --- периферійно для теми, хоча
Марценюк наводить мікрохвильову методику вимірювання обертання плазми
($\omega\approx1.16\times10^4$~рад/с на стелараторі).

\subsection{Кінетичні ефекти}

Москвитина \cite{Moskvitina2010} --- повноорбітне трасування $\alpha$-частинки в
полі ITER з гофруванням (реальні параметри, відтворювана схема; РК-5 в $\approx7$ разів
ефективніший за РК-4). Hoppe та ін.\ \cite{Hoppe2022} і Марчук \cite{Marchuk2025} ---
баланс електронів-утікачів
\begin{equation}\label{eq:RE}
  \frac{dn_{re}}{dt}=\gamma_{\rm D}+\frac{n_{re}}{t_{\rm av}}-\frac{n_{re}}{\tau_{re}},\qquad
  E_c=\frac{e^3n_e\ln\Lambda}{4\pi\varepsilon_0^2m_ec^2},
\end{equation}
з верифікацією STREAM на розряді JET \#77210. Тищенко \cite{Tyshchenko2021} описує
перенос енергії поперек поверхонь альфвеновими модами. Для моделі \emph{потоку}
кластер допоміжний: трасування частинок --- спосіб перевірки полів, а утікачі
важливі лише для сценаріїв старту і зриву.

\subsection{Інженерія, загальні й непрофільні роботи}

Зайцев та ін.\ \cite{Zaitsev2003,Zaitsev2001} --- МСЕ-механіка обмоток (корисні хіба
для геометрії котушок). Mazzucato \cite{Mazzucato2000} --- скейлінги $\tau_E$ і ліміт
Грінвальда (для вступу й оцінок порядків). Зимников \cite{Zimnikov2016} ---
популярні тези, науковим джерелом не є. Толочкевич та ін.\ \cite{Tolochkevych2015}
--- PIC-моделювання в незамагніченій плазмі (методичний приклад).
Аксенов \cite{Aksenov2015} містить розмірнісні помилки й суперечить
загальноприйнятій фізиці; \textbf{цитувати не рекомендується}.

%======================================================================
\section{Синтез: каркас моделі потоку плазми для проєкту}

На основі корпусу пропонується ієрархічна модель (у дужках --- джерела-основи):

\begin{enumerate}
  \item \textbf{Геометрія.} Рівновага (\ref{eq:GS}) з фіксованою або вільною межею;
        потокові координати $(\psi,\theta,\varphi)$
        \cite{Merlo2011,Seto2014}. Для JET --- D-подібний переріз, $\kappa\approx1.7$.
  \item \textbf{Транспорт 1.5D.} Рівняння (\ref{eq:n})--(\ref{eq:B}) для
        $n_e,n_i,T_e,T_i,\psi$ \cite{Fukuyama1992,Tyshchenko2021}.
  \item \textbf{Тороїдальний імпульс.} Додаткове рівняння для $L_\varphi=m_in_iRv_\varphi$
        з дифузійним і недифузійним потоками (\ref{eq:Gamma}), перезарядним
        гальмуванням (\ref{eq:mom}) і джерелом NBI \cite{Singh2004,Ida2001,Sigmar1987}.
  \item \textbf{Полоїдальний потік і $E_r$.} Неокласичний полоїдальний потік
        (у режимі Пфірша--Шлютера $U_\theta\approx-1.83\,\frac{T_i}{eB}\partial_r\ln T_i$
        \cite{Singh2004}) та $E_r$ з балансу сил (\ref{eq:Er}).
  \item \textbf{Зворотний зв'язок на транспорт.} Шир $E\times B$ пригнічує
        турбулентність при $\gamma_E\ge\gamma_T$ \cite{Singh2004,Wang2006} ---
        може бути введений як множник до $\chi_s$.
  \item \textbf{Граничні умови.} На сепаратрисі --- задані $n,T,v_\varphi$ або
        спрощена SOL-модель з умовами Бома (\ref{eq:Bohm}) \cite{AhoMantila2011}.
  \item \textbf{Валідація.} Порядки величин $E_r$, $v_\theta$, $v_\varphi$
        \cite{Ida2001}; профілі $T_{e,i}$ JET DTE1 \cite{Tyshchenko2021}; для
        перевірки розв'язувача рівноваги --- дані \cite{Merlo2011}.
\end{enumerate}

\noindent Рівень 2D (ядро + SOL в єдиній моделі) описано в \cite{Seto2014}, але для
студентського проєкту розумно почати з 1.5D і лише потім переходити до 2D.

%======================================================================
\section{Прогалини корпусу та рекомендації}

\begin{enumerate}
  \item \textbf{Немає робіт, що моделюють потоки на JET.} JET присутній лише як
        верифікаційний розряд \cite{Hoppe2022}, параметри в таблиці \cite{vanVugt2019}
        або розрахунок енергетичного балансу \cite{Tyshchenko2021}. Потрібні джерела
        з вимірювань обертання на JET (CXRS) і моделювання кодами JINTRAC/JETTO.
  \item \textbf{Немає неокласики для бананного режиму.} Усі теоретичні роботи про потоки
        --- для високої зіткненнєвості. Потрібні класичні огляди
        (Hinton--Hazeltine; Hirshman--Sigmar) та підручник рівня Wesson «Tokamaks».
  \item \textbf{Не отримано ключові класичні статті зі списку}: Муховатов--Шафранов
        (рівновага, 1971), Арцимович (1972), Furth (1975) і огляд De Tommasi
        (магнітне керування, 2018). Вони у відкритому доступі --- варто завантажити
        вручну через браузер.
  \item \textbf{ELAKPI (13 позицій)} --- зокрема магістерська Трегубова (моделювання
        магнітного підвісу) і тези Сливки (утримання плазми) --- доступні з мережі КПІ.
  \item \textbf{Вітчизняний сегмент корпусу} майже не стосується потоків плазми;
        для цієї теми опорними є англомовні джерела.
\end{enumerate}

%======================================================================
\section{Якість джерел і уточнення метаданих}

\paragraph{Невідповідності в назвах файлів.} У назвах файлів у теці \file{library}
у кількох випадках вказано не першого автора. Коректні дані наведено в бібліографії
нижче: C01 (перший автор Singh), C02 (Zou, Li, Lv), C03 (Sigmar), C04 (Wong),
C06 (Hoppe), C08 (Fable), C11 (van Vugt), C12 (Ida), C15 (Pigarov), C19 (Fukuyama),
C20 (W.\,X.~Wang). Запис CORE для C02 містить хибний DOI (J.\ Amer.\ Math.\ Soc., 1989),
узятий, найімовірніше, з першого посилання статті.

\paragraph{Якість тексту.} Скани 1987--1992 рр.\ (C03, C04, C05, C09, C13, C17) мають
пошкоджений OCR; у C20, G1, Н1, Н2 текстовий шар зламаний (кодування). Зміст цих
документів відновлено за зображеннями сторінок; формули варто звіряти з
транскрипціями Marker або з оригіналом. C14 містить лише оглядову частину
дисертації (статті I--V відсутні).

\paragraph{Надійність.} Н1 \cite{Aksenov2015} --- ненадійне; Х3 \cite{Zimnikov2016}
--- ненаукове; Н2 \cite{Duba2018} --- модель без обґрунтування параметрів.

%======================================================================
\begin{thebibliography}{99}
\small

\bibitem{Singh2004} Singh R., Rogister A., Kaw P. The effect of asymmetric gas puffing
  on toroidal flow in the edge of tokamak plasma~// Physics of Plasmas. -- 2004. --
  Vol.~11, No.~1. -- P.~129--139. -- DOI: 10.1063/1.1633558. [C01]

\bibitem{Zou2013} Zou W., Li F., Lv B. On a nonlocal problem for a confined plasma in
  a Tokamak~// Applications of Mathematics. -- 2013. -- Vol.~58, No.~6. --
  P.~609--642. [C02]

\bibitem{Sigmar1987} Sigmar D.\,J., Zanino R., Hsu C.\,T. Evolution of poloidal
  variation of impurity density and ambipolar potential in rotating tokamak plasma.
  Part~I~: Report PFC/RR-87-8. -- MIT Plasma Fusion Center, 1987. -- 16~p. [C03]

\bibitem{Wong1987} Wong K.\,L., Cheng C.\,Z. Neoclassical diffusion of heavy
  impurities in a rotating tokamak plasma~: Report PPPL-2470. -- Princeton Plasma
  Physics Laboratory, 1987. -- 13~p. [C04]

\bibitem{Ward1992} Ward D.\,J., Jardin S.\,C. The effects of plasma deformability on
  the feedback stabilization of axisymmetric modes in tokamak plasmas~: Report
  PPPL-2805. -- Princeton Plasma Physics Laboratory, 1992. -- 59~p. [C05]

\bibitem{Hoppe2022} Hoppe M., Ekmark I., Berger E., Fülöp T. Runaway electron
  generation during tokamak start-up~// Journal of Plasma Physics. -- 2022. --
  Vol.~88. -- 905880317. -- DOI: 10.1017/S002237782200054X. [C06]

\bibitem{Mazzucato2000} Mazzucato E. The importance of the toroidal magnetic field
  for the feasibility of a tokamak burning plasma experiment~: preprint. -- Princeton
  Plasma Physics Laboratory, 2000. -- 9~p. [C07]

\bibitem{Fable2022} Fable E., Janky F., Treutterer W. et al. The modeling of a
  tokamak plasma discharge, from first principles to a flight simulator~// Plasma
  Physics and Controlled Fusion. -- 2022. -- Vol.~64. -- 044002. --
  DOI: 10.1088/1361-6587/ac466b. [C08]

\bibitem{Park1989} Park H.\,K., Goldston R.\,J., Taylor G. Study of electron density
  and pressure profile behavior in a tokamak plasma~: Report PPPL-2660. -- Princeton
  Plasma Physics Laboratory, 1989. -- 13~p. [C09]

\bibitem{Seto2014} Seto H., Fukuyama A. Formulation of two-dimensional transport
  modeling in tokamak plasmas~// Plasma and Fusion Research. -- 2014. -- Vol.~9. --
  1403002. -- DOI: 10.1585/pfr.9.1403002. [C10]

\bibitem{vanVugt2019} van Vugt D.\,C., Huijsmans G.\,T.\,A., Hoelzl M., Loarte A.
  Kinetic modeling of ELM-induced tungsten transport in a tokamak plasma~// Physics of
  Plasmas. -- 2019. -- Vol.~26, No.~4. -- 042508. -- DOI: 10.1063/1.5092319. [C11]

\bibitem{Ida2001} Ida K., CHS Group, LHD Group, JFT-2M Group. Structure of the radial
  electric field and toroidal/poloidal flow in high temperature toroidal plasma~//
  J.~Plasma Fusion Res. SERIES. -- 2001. -- Vol.~4. -- P.~36--42. [C12]

\bibitem{Hanada1990} Hanada K., Tanaka H., Iida M. et al. Sawtooth stabilization by
  localized electron cyclotron heating in a tokamak plasma~: Report PPPL-2725. --
  Princeton Plasma Physics Laboratory, 1990. -- 16~p. [C13]

\bibitem{AhoMantila2011} Aho-Mantila L. Divertor plasma conditions and their effect on
  carbon migration in the ASDEX Upgrade tokamak~: doctoral dissertation. -- Espoo~:
  Aalto University; VTT Publications 773, 2011. [C14]

\bibitem{Pigarov2006} Pigarov A.\,Yu., Smirnov R.\,D., Krasheninnikov S.\,I.,
  Rognlien T.\,D., Rozenberg M. Transport of dust particles in tokamak devices~:
  preprint UCRL-JRNL-221928. -- LLNL, 2006. -- 17~p. [C15]

\bibitem{Merlo2011} Merlo P. Plasma boundary reconstruction and shape control of
  tokamak discharges in RFX-mod~: tesi di laurea magistrale. -- Padova~: Università
  degli Studi di Padova, 2011. -- 170~p. [C16]

\bibitem{Hulse1992} Hulse R.\,A. Modeling of impurity transport in the core plasma~:
  Report PPPL-CFP-2776. -- Princeton Plasma Physics Laboratory, 1992. -- 28~p. [C17]

\bibitem{Jardin2000} Jardin S.\,C., Kaye S., Menard J., Kessel C., Glasser A.\,H.
  Tokamak Simulation Code modeling of NSTX~: Report PPPL-3472. -- Princeton Plasma
  Physics Laboratory, 2000. -- 7~p. [C18]

\bibitem{Fukuyama1992} Fukuyama A., Kasai T., Furutani Y. Transport simulasion [sic]
  in a burning tokamak plasma~// Memoirs of the Faculty of Engineering, Okayama
  University. -- 1992. -- Vol.~26, No.~2. -- P.~93--109. [C19]

\bibitem{Wang2006} Wang W.\,X., Lin Z., Tang W.\,M., Lee W.\,W. et al. Gyrokinetic
  simulation of global turbulent transport properties in tokamak experiments~: Report
  PPPL-4142. -- Princeton Plasma Physics Laboratory, 2006. -- 34~p. [C20]

\bibitem{Martsenyuk2025} Марценюк Ю.\,П. Характеристики газорозрядної плазми та
  розрядів різних типів у термоядерних дослідженнях~: дис. \ldots\ д-ра філософії~:
  104. -- Харків~: ННЦ ХФТІ НАН України, 2025. -- 150~с. [G1]

\bibitem{Marchuk2025} Марчук Ю.\,М. Моделювання електронів-утікачів, парних зіткнень
  та процесів у нижньому гібридному резонансі в плазмі з магнітним полем~: дис.
  \ldots\ д-ра філософії~: 104. -- Харків~: ІФП ННЦ ХФТІ НАН України, 2025. --
  147~с. [G2]

\bibitem{Romashchenko2021} Ромащенко О.\,В. Динаміка та фазові стани макрочастинок у
  пучково-плазмових системах~: автореф. дис. \ldots\ д-ра фіз.-мат. наук~: 01.04.08.
  -- Харків~: ХНУ ім. В.\,Н.~Каразіна, 2021. -- 36~с. [G3]

\bibitem{Tyshchenko2021} Тищенко М.\,Г. Поширення альфвенових хвиль та перенесення
  енергії поперек магнітних поверхонь у тороїдальній плазмі~: дис. \ldots\ канд.
  фіз.-мат. наук~: 01.04.08. -- Харків, 2021. -- 140~с. [G4]

\bibitem{Tolochkevych2015} Толочкевич Ю.\,М., Анісімов І.\,О., Літошенко Т.\,Є.
  Динаміка заряджених згустків у збудженому ними в плазмі кільватерному полі~//
  Український фізичний журнал. -- 2015. -- Т.~60, №~1. -- С.~16--22. [G5]

\bibitem{Aksenov2015} Аксенов В.\,В. Векторный потенциал, эффект Ааронова--Бома и
  электродинамика в токамаке~// Праці Міжнародного геометричного центру. -- 2015. --
  Т.~8, №~3--4. -- С.~31--39. [Н1]

\bibitem{Duba2018} Дуба І.\,Е. Сучасна цифрова багатозв'язна система управління
  основним контуром -- положення плазмового шнура -- установки термоядерного синтезу
  ТОКАМАК~// Вчені записки ТНУ ім. В.\,І.~Вернадського. Сер.: Технічні науки. --
  2018. -- Т.~29(68), №~4(1). -- С.~132--137. [Н2]

\bibitem{Moskvitina2010} Москвитина Ю.\,К. Численное моделирование влияния
  гофрировки на движение заряженной частицы в токамаке ITER~// Вісник ХНУ ім.
  В.\,Н.~Каразіна. Сер.: Математичне моделювання. -- 2010. -- №~925, вип.~14. --
  С.~147--154. [Н3]

\bibitem{Zaitsev2003} Зайцев Б.\,Ф., Асаёнок А.\,В. Напряженное состояние в винтовой
  обмотке тороидальной магнитной системы электрофизической установки под действием
  температуры~// Вестник НТУ «ХПИ». Динамика и прочность машин. -- 2003. -- №~12,
  т.~1. -- С.~64--71. [Х1]

\bibitem{Zaitsev2001} Зайцев Б.\,Ф., Шульженко Н.\,Г. Управление напряженным
  состоянием плоской многослойной катушки тороидального магнитного поля~// Вестник
  НТУ «ХПИ». Динамика и прочность машин. -- 2001. -- №~25. -- С.~86--90. [Х2]

\bibitem{Zimnikov2016} Зимников О.\,О., Дьяконенко Н.\,Л., Макогон О.\,А.
  Термоядерный синтез -- энергетика будущего~// Актуальні проблеми фізики та їх
  інформаційне забезпечення~: тези доп. -- Харків~: НТУ «ХПІ», 2016. -- С.~82--83. [Х3]

\bibitem{Belyaeva2016} Беляева А.\,И., Галуза А.\,А., Коленов И.\,В., Савченко А.\,А.
  Роль рекристаллизации вольфрама в формировании шероховатости его поверхности под
  влиянием последовательного воздействия нейтронов и распыления~// Металлофизика и
  новейшие технологии. -- 2016. -- Т.~38, №~8. -- С.~1077--1102. [Х4]

\end{thebibliography}

\end{document}
